% GENERATED FILE. Edit canonical lessons, then run npm run generate:books.
% canonical-source-hash: fnv1a64:07f641b2ca874ed5
% canonical-lessons: JA-C430-ikari, JA-C430-nonoshiru, JA-C430-niramu, JA-C430-iya, JA-C430-abareru

\chapter{Anger and Dislike}
\label{ch:ja-anger-and-dislike}
\hlchaptermodality{430}

\begin{chapteropening}
\textbf{By the end of this chapter:} \emph{I can talk about anger and things I dislike.}

Complete the last lesson of chapter 430: I can talk about anger and things I dislike.
\end{chapteropening}

\section[\texorpdfstring{\ja{いかり} (ikari)}{いかり (ikari)}]{\ja{いかり} (ikari) — anger; rage}

\label{lesson:JA-C430-ikari}



\begin{quote}
Before the new one: say the Japanese for hard to bear, then the Japanese for to comfort.
\end{quote}

\subsection*{You'll want to know: \ja{いかり}}
\textbf{\ja{いかり}} — \emph{ikari} — \textquotedblleft{}anger; rage\textquotedblright{}.

\begin{quote}
\textbf{\ja{ちちの} \ja{いかりは} \ja{よく} \ja{わかります。}} — \emph{chichi no ikari wa yoku wakarimasu} — \textquotedblleft{}I understand my father's anger well.\textquotedblright{}
\end{quote}

\begin{cousinweb}
\textbf{\ja{いかり}} (\emph{ikari}) is a native Japanese word (wago), inherited rather than borrowed from Chinese; its kanji spelling was attached to it later (JMdict).
\end{cousinweb}

\begin{grammarlens}[title={What you've built}]
One more word for describing what you lived and felt.
\end{grammarlens}

\subsection*{Your turn}
\begin{itemize}
\raggedright
  \item \emph{Say it:} \emph{ikari}
  \item \emph{Say it:} \emph{ikari}, once more
  \item \emph{Say it:} \emph{chichi no ikari wa yoku wakarimasu} — I understand my father's anger well.
  \item \emph{Recall:} say the Japanese for hard to bear, then the Japanese for to comfort, then say \emph{ikari} again
\end{itemize}

\begin{tcolorbox}[breakable,colback=teal!4,colframe=teal!35!black,title={Before you move on}]
What does \ja{いかり} mean? (\textquotedblleft{}Anger; rage\textquotedblright{}.) Say it once more.
\end{tcolorbox}
\section[\texorpdfstring{\ja{ののしる} (nonoshiru)}{ののしる (nonoshiru)}]{\ja{ののしる} (nonoshiru) — to curse at; to shout abuse at}

\label{lesson:JA-C430-nonoshiru}



\begin{quote}
Before the new one: say the Japanese for to comfort, then the Japanese for anger.
\end{quote}

\subsection*{You'll want to know: \ja{ののしる}}
\textbf{\ja{ののしる}} — \emph{nonoshiru} — \textquotedblleft{}to curse at; to shout abuse at\textquotedblright{}.

\begin{quote}
\textbf{\ja{ともだちを} \ja{ののしらないで} \ja{ください。}} — \emph{tomodachi o nonoshiranaide kudasai} — \textquotedblleft{}Please do not swear at your friends.\textquotedblright{}
\end{quote}

Polite \textbf{\ja{ののしります}}, \emph{nonoshirimasu}; a godan verb. Past (polite): \textbf{\ja{かれが} \ja{てんいんを} \ja{ののしりました。}}, \emph{kare ga ten'in o nonoshirimashita}.

\begin{cousinweb}
\textbf{\ja{ののしる}} (\emph{nonoshiru}) is a native Japanese word (wago), inherited rather than borrowed from Chinese; its kanji spelling was attached to it later (JMdict).
\end{cousinweb}

\begin{grammarlens}[title={What you've built}]
One more word for describing what you lived and felt.
\end{grammarlens}

\subsection*{Your turn}
\begin{itemize}
\raggedright
  \item \emph{Say it:} \emph{nonoshiru}
  \item \emph{Say it:} \emph{nonoshiru}, once more
  \item \emph{Say it:} \emph{tomodachi o nonoshiranaide kudasai} — Please do not swear at your friends.
  \item \emph{Recall:} say the Japanese for to comfort, then the Japanese for anger, then say \emph{nonoshiru} again
\end{itemize}

\begin{tcolorbox}[breakable,colback=teal!4,colframe=teal!35!black,title={Before you move on}]
What does \ja{ののしる} mean? (\textquotedblleft{}To curse at; to shout abuse at\textquotedblright{}.) Say it once more.
\end{tcolorbox}
\section[\texorpdfstring{\ja{にらむ} (niramu)}{にらむ (niramu)}]{\ja{にらむ} (niramu) — to glare at}

\label{lesson:JA-C430-niramu}



\begin{quote}
Before the new one: say the Japanese for anger, then the Japanese for to curse at.
\end{quote}

\subsection*{You'll want to know: \ja{にらむ}}
\textbf{\ja{にらむ}} — \emph{niramu} — \textquotedblleft{}to glare at\textquotedblright{}.

\begin{quote}
\textbf{\ja{ねこが} \ja{いぬを} \ja{にらんでいます。}} — \emph{neko ga inu o nirande imasu} — \textquotedblleft{}The cat is glaring at the dog.\textquotedblright{}
\end{quote}

Polite \textbf{\ja{にらみます}}, \emph{niramimasu}; a godan verb. Past (polite): \textbf{\ja{あにが} \ja{おとうとを} \ja{にらみました。}}, \emph{ani ga otōto o niramimashita}.

\begin{cousinweb}
\textbf{\ja{にらむ}} (\emph{niramu}) is a native Japanese word (wago), inherited rather than borrowed from Chinese; its kanji spelling was attached to it later (JMdict).
\end{cousinweb}

\begin{grammarlens}[title={What you've built}]
One more word for describing what you lived and felt.
\end{grammarlens}

\subsection*{Your turn}
\begin{itemize}
\raggedright
  \item \emph{Say it:} \emph{niramu}
  \item \emph{Say it:} \emph{niramu}, once more
  \item \emph{Say it:} \emph{neko ga inu o nirande imasu} — The cat is glaring at the dog.
  \item \emph{Recall:} say the Japanese for anger, then the Japanese for to curse at, then say \emph{niramu} again
\end{itemize}

\begin{tcolorbox}[breakable,colback=teal!4,colframe=teal!35!black,title={Before you move on}]
What does \ja{にらむ} mean? (\textquotedblleft{}To glare at\textquotedblright{}.) Say it once more.
\end{tcolorbox}
\section[\texorpdfstring{\ja{いや} (iya)}{いや (iya)}]{\ja{いや} (iya) — unpleasant; unwanted}

\label{lesson:JA-C430-iya}



\begin{quote}
Before the new one: say the Japanese for to curse at, then the Japanese for to glare at.
\end{quote}

\subsection*{You'll want to know: \ja{いや}}
\textbf{\ja{いや}} — \emph{iya} — \textquotedblleft{}unpleasant; unwanted\textquotedblright{}.

\begin{quote}
\textbf{\ja{さむい} \ja{あさは} \ja{いやです。}} — \emph{samui asa wa iya desu} — \textquotedblleft{}I do not like cold mornings.\textquotedblright{}
\end{quote}

A na-adjective: \ja{な} before a noun.

\begin{cousinweb}
\textbf{\ja{いや}} (\emph{iya}) is a native Japanese word (wago), inherited rather than borrowed from Chinese; its kanji spelling was attached to it later (JMdict).
\end{cousinweb}

\begin{grammarlens}[title={What you've built}]
One more word for describing what you lived and felt.
\end{grammarlens}

\subsection*{Your turn}
\begin{itemize}
\raggedright
  \item \emph{Say it:} \emph{iya}
  \item \emph{Say it:} \emph{iya}, once more
  \item \emph{Say it:} \emph{samui asa wa iya desu} — I do not like cold mornings.
  \item \emph{Recall:} say the Japanese for to curse at, then the Japanese for to glare at, then say \emph{iya} again
\end{itemize}

\begin{tcolorbox}[breakable,colback=teal!4,colframe=teal!35!black,title={Before you move on}]
What does \ja{いや} mean? (\textquotedblleft{}Unpleasant; unwanted\textquotedblright{}.) Say it once more.
\end{tcolorbox}
\section[\texorpdfstring{\ja{あばれる} (abareru)}{あばれる (abareru)}]{\ja{あばれる} (abareru) — to run wild; to act violently}

\label{lesson:JA-C430-abareru}



\begin{quote}
Before the new one: say the Japanese for to glare at, then the Japanese for unpleasant.
\end{quote}

\subsection*{You'll want to know: \ja{あばれる}}
\textbf{\ja{あばれる}} — \emph{abareru} — \textquotedblleft{}to run wild; to act violently\textquotedblright{}.

\begin{quote}
\textbf{\ja{うまが} \ja{あばれて、おとこが} \ja{けがを} \ja{しました。}} — \emph{uma ga abarete, otoko ga kega o shimashita} — \textquotedblleft{}A horse ran wild and a man was hurt.\textquotedblright{}
\end{quote}

Polite \textbf{\ja{あばれます}}, \emph{abaremasu}; an ichidan verb. Past (polite): \textbf{\ja{うしが} \ja{あばれました。}}, \emph{ushi ga abaremashita}.

\begin{cousinweb}
\textbf{\ja{あばれる}} (\emph{abareru}) is a native Japanese word (wago), inherited rather than borrowed from Chinese; its kanji spelling was attached to it later (JMdict).
\end{cousinweb}

\begin{grammarlens}[title={What you've built}]
One more word for describing what you lived and felt.
\end{grammarlens}

\subsection*{Your turn}
\begin{itemize}
\raggedright
  \item \emph{Say it:} \emph{abareru}
  \item \emph{Say it:} \emph{abareru}, once more
  \item \emph{Say it:} \emph{uma ga abarete, otoko ga kega o shimashita} — A horse ran wild and a man was hurt.
  \item \emph{Recall:} say the Japanese for to glare at, then the Japanese for unpleasant, then say \emph{abareru} again
\end{itemize}

\begin{tcolorbox}[breakable,colback=teal!4,colframe=teal!35!black,title={Before you move on}]
What does \ja{あばれる} mean? (\textquotedblleft{}To run wild; to act violently\textquotedblright{}.) Say it once more.
\end{tcolorbox}
